\chapter{连续极限与 Symanzik 改进方案}

对规范场与费米子自由度的简单威尔逊作用量作泰勒展开，显示出高阶维算符的存在。由于这些算符被格距的幂次所压低，在重整化群的语言中它们属于无关算符，即当 $a \to 0$ 时它们在不动点处消失。表达同一事实的另一种方式如下。以 $\pi/a$ 为硬截断的格点理论可以被视为一种有效理论 [49]。把动量从 $\infty \to \pi/a$ 积分掉会产生有效相互作用。因此，一般情况下我们可以把作用量写成
\begin{equation}
S_C = S_L(\pi/a) + \sum_n \sum_i a^n C_i^n A_i^n(\pi/a) \,,
\label{eq:8.1}
\end{equation}
其中 $S_C$（$S_L$）是连续（格点）作用量，第一个求和遍历给定维数的所有算符，第二个求和遍历所有大于四的维数。类似地，用来探测物理的算符可以写成
\begin{equation}
O_C = O_L(\pi/a) + \sum_n \sum_i a^n D_i^n O_i^n(\pi/a) \,,
\label{eq:8.2}
\end{equation}
在此情形下，量子修正还会诱导与低维算符的混合。系数 $C_i^n$ 和 $D_i^n$ 都是耦合常数 $\alpha_s$ 的函数。附带提一句，在有限 $a$ 下保持规范、洛伦兹和手征对称性的一个重要原因是，它们极大地限制了给定维数下可能算符的集合。

现在考虑某个给定算符的期望值
\begin{equation}
\begin{split}
\langle O_C \rangle &= \int dU\, O_C e^{-S_C} \\
&= \int dU \left( O_L + \sum_n \sum_i a^n D_i^n O_i^n \right)
   e^{-S_L - \sum_m \sum_i a^m C_i^m A_i^m} \\
&= \int dU \left( O_L + \sum_n \sum_i a^n D_i^n O_i^n \right)
   \left( 1 - \sum_m \sum_i a^m C_i^m A_i^m + \cdots \right) e^{-S_L} \\
&= \int dU\, O_L e^{-S_L}
   + \int dU\, O_L \left( \sum_m \sum_i a^m C_i^m A_i^m + \cdots \right) e^{-S_L} \\
&\qquad + \int dU \left( \sum_n \sum_i a^n D_i^n O_i^n \right)
   \left( 1 - \sum_m \sum_i a^m C_i^m A_i^m + \cdots \right) e^{-S_L} \,.
\end{split}
\label{eq:8.3}
\end{equation}
注意，出现在作用量中的算符 $A_i^n$ 要对所有时空点求和。这会在式~\ref{eq:8.3} 中导致接触项；为简洁起见，我将忽略这些接触项，因为它们需要对算符给出适当的定义，但并不改变结论。

相信 LQCD 在有限 $a$ 下通过模拟能够提供可靠结果的基础是，格点赝象的贡献（最后一个表达式中的后两项）很小、可以被估计并因而被移除。在这一点上，我只就期望值的计算作几点一般性的陈述，这些将在后面的讲义中详加阐述。
\begin{itemize}
\item 式~\ref{eq:8.3} 中最后两项的贡献在 $a \to 0$ 极限下消失，除非存在与低维算符的混合。在这种情况下，混合系数必须被非常精确地确定（这需要非微扰方法），否则 $a \to 0$ 极限是发散的。
\item 计算必须在这样的 $a$ 值下进行：其中按 $a$ 的连续幂次组织的无关项贡献至少按几何级数递减。
\item 为获得连续极限结果，模拟需在足够多的 $a$ 值下进行，以便能可靠地外推到 $a = 0$。进行计算所用格距 $a$ 的「精细程度」以及所需的取值点数，取决于修正的大小以及我们对用于外推的函数形式能够确定到何种程度。
\item 为了把结果改进到 $a$ 的任意给定阶，必须把作用量和算符都改进到相同的阶。这种基于去除按 $a$ 的幂级数展开来组织的格点赝象的方法，称为 Symanzik 改进方案。
\end{itemize}

式~\ref{eq:8.3} 清楚地表明，构造 LQCD 有相当大的灵活性。我们可以自由地以合理的强度加入任何无关算符，并且在连续极限下仍然恢复 QCD。作用量的各种构造之间的唯一差别在于趋近连续极限的方式。我们想做的是改进作用量和算符，以便在尽可能粗的格点上得到连续极限结果。另一方面，改进要求给作用量和算符都加上无关项。这会增加计算的复杂性和模拟时间。因此，最小化离散化误差的策略是一种折衷——在作用量和算符的简单性（以模拟成本、代码实现和分析来度量）与在某个给定的格距值（例如 $a \approx 0.1$ 费米）处评估的离散化误差的减小之间进行优化。剩余的误差（可能很小）随后可以通过把从少数几个 $a$ 值下模拟得到的结果外推到 $a = 0$ 来消除。

LQCD 模拟所需的浮点运算次数在淬火近似下大致按 $L^6$ 标度，而在具有动力学费米子时按 $L^{10}$ 标度。当前的 QCD 格点模拟所用的格距 $a$ 在 $2\,\mathrm{GeV} \leq a^{-1} \leq 5\,\mathrm{GeV}$ 范围内变化。对于这些参数值，在许多可观测量中发现威尔逊作用量中 $O(a)$、$O(ma)$、$O(pa)$、$O(a^2)$、\ldots\ 各项引起的修正是很大的。因此，改进格点表述（作用量和算符）的重要性不言而喻，对于完整 QCD 尤其如此。因此，寻找这样的作用量是当下的热门课题。在讨论这些之前，让我先总结一下人们希望从这样一个改进方案中获得什么，以及用什么标准来判断它。
\begin{enumerate}
\item 改进的标度行为，即离散化误差被减小。因此，在给定精度下可以在更粗的格点（更小的 $L = N_S$）上进行计算。即使 $L$ 减小 2 倍，也对应计算机时间上 $2^{10} \approx 1000$ 倍的缩减！

\item 更好地恢复转动对称性和内部对称性，例如手征、交错味对称性等。

\item 取连续极限所沿的轨迹不应过于靠近作为格点离散化赝象的那些额外相变（关于格点理论中相变的讨论见第 15 节）。此类相变的问题在于，在它们附近，所期望的连续标度行为可能会对某些可观测量发生改变，即式~\ref{eq:8.3} 中的赝象项在 $a$ 上可能不是几何级数收敛的，因而可能很大。

\item 格点上的拓扑研究对短距离涨落尤为敏感，这些涨落是离散化规范作用量的赝象。对于威尔逊规范作用量，人们发现用 Lüscher 的几何方法 [50] 计算得到的拓扑荷，在瞬子核心半径 $\rho_c a \approx 0.7$ 处从 $0 \to 1$ 跳跃，而这类瞬子的作用量 $S_{\rho_c}$ 明显小于经典值 $S^{\mathrm{cl}} = 8\pi^2/\beta$。对于这种 $S < S^{\mathrm{cl}}$ 的小瞬子，熵因子压倒了玻尔兹曼抑制，导致连续极限下拓扑磁化率发散。这些短距离赝象被称为位错，当 $\rho \approx \rho_c$ 时，如何区分物理与非物理的部分，并非事后（a posteriori）显而易见。因此，从一个排除此类赝象的格点表述出发非常重要，这样的表述中只存在物理的瞬子，即那些平均核心大小对 $g$ 具有正确标度行为的瞬子。人们发现，改进作用量会抑制位错，即改进规范作用量会增大 $S_{\rho_c}$ [51]，并且对拓扑荷给出更好的定义 [52]。

\item 随着越来越多的无关算符被加入作用量，格点微扰理论变得越来越复杂。尽管随着用于确定重整化常数的非微扰方法的出现，对格点微扰理论的依赖已经减少，但微扰分析仍不失为有价值的指导和检验，因此我们希望保留这种可能性。对于改进作用量，人们期望的是：$\alpha_s$ 在格点和连续理论中的值，以及某个给定可观测量按 $\alpha_s$ 展开的级数系数，彼此更为接近。在这种情况下，两种理论之间的匹配要可靠得多。目前，格点理论与连续理论之间的 1 圈匹配是较大的不确定来源之一。

\item Adler–Bell–Jackiw（ABJ）反常：在连续理论中，味单态轴矢流具有反常。这作为 UV 正则化方案的一部分而出现。在 Pauli–Villars 方案中，引入一个重质量作为截断，这会破坏 $\gamma_5$ 不变性；而在维数正则化中，$\gamma_5$ 在离开 4 维时没有良好的定义。如下所示，朴素费米子保持了 $\gamma_5$ 不变性，但以味加倍为代价。加倍费米子的一个后果是成对相消，从而没有净反常。一个可行的格点理论应当重现 ABJ 反常。
\end{enumerate}

在讨论这样的「改进」作用量之前，重新审视朴素狄拉克作用量的构造以凸显其问题，是有益的。
